% =============================================================================
% figures/figura_diagrama_fases_agua.tex
% Diagrama de fases esquematico del agua — generado con entorno picture
% Se incluye en el capitulo con: % =============================================================================
% figures/figura_diagrama_fases_agua.tex
% Diagrama de fases esquematico del agua — generado con entorno picture
% Se incluye en el capitulo con: % =============================================================================
% figures/figura_diagrama_fases_agua.tex
% Diagrama de fases esquematico del agua — generado con entorno picture
% Se incluye en el capitulo con: % =============================================================================
% figures/figura_diagrama_fases_agua.tex
% Diagrama de fases esquematico del agua — generado con entorno picture
% Se incluye en el capitulo con: \input{figures/figura_diagrama_fases_agua}
% =============================================================================

\begin{figure}[htbp]
\centering
\setlength{\unitlength}{1.2cm}
\begin{picture}(9, 7)

    % ---- Ejes ---------------------------------------------------------------
    \put(1.0, 0.5){\vector(1, 0){7.5}}   % Eje X (Temperatura)
    \put(1.0, 0.5){\vector(0, 1){6.0}}   % Eje Y (Presion)

    % Etiquetas de ejes
    \put(8.7, 0.25){\small $T$ / \si{\celsius}}
    \put(0.0, 6.6){\small $p$ / \si{\bar}}

    % ---- Lineas de las curvas de coexistencia --------------------------------

    % Curva solido-liquido (ligeramente inclinada a la izquierda — anomalia del agua)
    \put(3.5, 1.7){\line(-1, 4){0.9}}     % Pendiente negativa (dP/dT < 0)
    % Etiqueta
    \put(2.0, 4.5){\small\itshape Curva de}
    \put(2.0, 4.2){\small\itshape fusi\'on}

    % Curva liquido-vapor (desde punto triple hasta punto critico)
    \qbezier(3.5, 1.7)(5.5, 2.5)(7.5, 5.5)
    \put(5.8, 3.0){\small\itshape Curva de}
    \put(5.8, 2.7){\small\itshape vaporizaci\'on}

    % Curva solido-vapor (desde 0 hasta punto triple)
    \qbezier(1.2, 0.7)(2.0, 0.9)(3.5, 1.7)
    \put(1.5, 0.9){\small\itshape Curva de}
    \put(1.5, 0.65){\small\itshape sublimaci\'on}

    % ---- Regions (etiquetas de fases) ---------------------------------------
    \put(1.8, 3.5){\large S\'OLIDO}
    \put(4.2, 2.2){\large L\'IQUIDO}
    \put(4.8, 1.0){\large VAPOR}

    % ---- Punto triple -------------------------------------------------------
    \put(3.5, 1.7){\circle*{0.12}}
    \put(3.2, 1.45){\small Punto triple}
    \put(3.2, 1.22){\small ($0.01$\,\si{\celsius}, $6.12\times10^{-3}$\,\si{\bar})}

    % ---- Punto critico ------------------------------------------------------
    \put(7.5, 5.5){\circle*{0.12}}
    \put(7.0, 5.8){\small Punto cr\'itico}
    \put(7.0, 5.6){\small ($374$\,\si{\celsius}, $221$\,\si{\bar})}

    % ---- Lineas de referencia (1 bar) ---------------------------------------
    \put(1.0, 2.0){\line(1, 0){7}}        % Linea horizontal p = 1 bar
    \put(8.1, 1.95){\tiny $p = 1$\,bar}

    % ---- Escala de presion --------------------------------------------------
    \put(0.6, 2.0){\small $1$}
    \put(0.3, 1.7){\small $6.1\times10^{-3}$}
    \put(0.6, 5.5){\small $221$}

    % ---- Escala de temperatura ----------------------------------------------
    \put(2.9, 0.25){\small $0.01$}
    \put(7.3, 0.25){\small $374$}

\end{picture}
\caption{Diagrama de fases esquem\'atico del agua (\ce{H2O}).
         Se muestran las tres curvas de coexistencia (s\'olido-l\'iquido,
         l\'iquido-vapor y s\'olido-vapor), el punto triple
         (\SI{0.01}{\celsius}, \SI{6.12e-3}{\bar}) y el punto cr\'itico
         (\SI{374}{\celsius}, \SI{221}{\bar}).
         N\'otese la pendiente \textit{negativa} de la curva de fusi\'on,
         caracter\'istica an\'omala del agua debida a que el hielo es menos
         denso que el agua l\'iquida.}
\label{fig:diagrama_fases_agua}
\end{figure}

% =============================================================================
% FIN DE figura_diagrama_fases_agua.tex
% =============================================================================

% =============================================================================

\begin{figure}[htbp]
\centering
\setlength{\unitlength}{1.2cm}
\begin{picture}(9, 7)

    % ---- Ejes ---------------------------------------------------------------
    \put(1.0, 0.5){\vector(1, 0){7.5}}   % Eje X (Temperatura)
    \put(1.0, 0.5){\vector(0, 1){6.0}}   % Eje Y (Presion)

    % Etiquetas de ejes
    \put(8.7, 0.25){\small $T$ / \si{\celsius}}
    \put(0.0, 6.6){\small $p$ / \si{\bar}}

    % ---- Lineas de las curvas de coexistencia --------------------------------

    % Curva solido-liquido (ligeramente inclinada a la izquierda — anomalia del agua)
    \put(3.5, 1.7){\line(-1, 4){0.9}}     % Pendiente negativa (dP/dT < 0)
    % Etiqueta
    \put(2.0, 4.5){\small\itshape Curva de}
    \put(2.0, 4.2){\small\itshape fusi\'on}

    % Curva liquido-vapor (desde punto triple hasta punto critico)
    \qbezier(3.5, 1.7)(5.5, 2.5)(7.5, 5.5)
    \put(5.8, 3.0){\small\itshape Curva de}
    \put(5.8, 2.7){\small\itshape vaporizaci\'on}

    % Curva solido-vapor (desde 0 hasta punto triple)
    \qbezier(1.2, 0.7)(2.0, 0.9)(3.5, 1.7)
    \put(1.5, 0.9){\small\itshape Curva de}
    \put(1.5, 0.65){\small\itshape sublimaci\'on}

    % ---- Regions (etiquetas de fases) ---------------------------------------
    \put(1.8, 3.5){\large S\'OLIDO}
    \put(4.2, 2.2){\large L\'IQUIDO}
    \put(4.8, 1.0){\large VAPOR}

    % ---- Punto triple -------------------------------------------------------
    \put(3.5, 1.7){\circle*{0.12}}
    \put(3.2, 1.45){\small Punto triple}
    \put(3.2, 1.22){\small ($0.01$\,\si{\celsius}, $6.12\times10^{-3}$\,\si{\bar})}

    % ---- Punto critico ------------------------------------------------------
    \put(7.5, 5.5){\circle*{0.12}}
    \put(7.0, 5.8){\small Punto cr\'itico}
    \put(7.0, 5.6){\small ($374$\,\si{\celsius}, $221$\,\si{\bar})}

    % ---- Lineas de referencia (1 bar) ---------------------------------------
    \put(1.0, 2.0){\line(1, 0){7}}        % Linea horizontal p = 1 bar
    \put(8.1, 1.95){\tiny $p = 1$\,bar}

    % ---- Escala de presion --------------------------------------------------
    \put(0.6, 2.0){\small $1$}
    \put(0.3, 1.7){\small $6.1\times10^{-3}$}
    \put(0.6, 5.5){\small $221$}

    % ---- Escala de temperatura ----------------------------------------------
    \put(2.9, 0.25){\small $0.01$}
    \put(7.3, 0.25){\small $374$}

\end{picture}
\caption{Diagrama de fases esquem\'atico del agua (\ce{H2O}).
         Se muestran las tres curvas de coexistencia (s\'olido-l\'iquido,
         l\'iquido-vapor y s\'olido-vapor), el punto triple
         (\SI{0.01}{\celsius}, \SI{6.12e-3}{\bar}) y el punto cr\'itico
         (\SI{374}{\celsius}, \SI{221}{\bar}).
         N\'otese la pendiente \textit{negativa} de la curva de fusi\'on,
         caracter\'istica an\'omala del agua debida a que el hielo es menos
         denso que el agua l\'iquida.}
\label{fig:diagrama_fases_agua}
\end{figure}

% =============================================================================
% FIN DE figura_diagrama_fases_agua.tex
% =============================================================================

% =============================================================================

\begin{figure}[htbp]
\centering
\setlength{\unitlength}{1.2cm}
\begin{picture}(9, 7)

    % ---- Ejes ---------------------------------------------------------------
    \put(1.0, 0.5){\vector(1, 0){7.5}}   % Eje X (Temperatura)
    \put(1.0, 0.5){\vector(0, 1){6.0}}   % Eje Y (Presion)

    % Etiquetas de ejes
    \put(8.7, 0.25){\small $T$ / \si{\celsius}}
    \put(0.0, 6.6){\small $p$ / \si{\bar}}

    % ---- Lineas de las curvas de coexistencia --------------------------------

    % Curva solido-liquido (ligeramente inclinada a la izquierda — anomalia del agua)
    \put(3.5, 1.7){\line(-1, 4){0.9}}     % Pendiente negativa (dP/dT < 0)
    % Etiqueta
    \put(2.0, 4.5){\small\itshape Curva de}
    \put(2.0, 4.2){\small\itshape fusi\'on}

    % Curva liquido-vapor (desde punto triple hasta punto critico)
    \qbezier(3.5, 1.7)(5.5, 2.5)(7.5, 5.5)
    \put(5.8, 3.0){\small\itshape Curva de}
    \put(5.8, 2.7){\small\itshape vaporizaci\'on}

    % Curva solido-vapor (desde 0 hasta punto triple)
    \qbezier(1.2, 0.7)(2.0, 0.9)(3.5, 1.7)
    \put(1.5, 0.9){\small\itshape Curva de}
    \put(1.5, 0.65){\small\itshape sublimaci\'on}

    % ---- Regions (etiquetas de fases) ---------------------------------------
    \put(1.8, 3.5){\large S\'OLIDO}
    \put(4.2, 2.2){\large L\'IQUIDO}
    \put(4.8, 1.0){\large VAPOR}

    % ---- Punto triple -------------------------------------------------------
    \put(3.5, 1.7){\circle*{0.12}}
    \put(3.2, 1.45){\small Punto triple}
    \put(3.2, 1.22){\small ($0.01$\,\si{\celsius}, $6.12\times10^{-3}$\,\si{\bar})}

    % ---- Punto critico ------------------------------------------------------
    \put(7.5, 5.5){\circle*{0.12}}
    \put(7.0, 5.8){\small Punto cr\'itico}
    \put(7.0, 5.6){\small ($374$\,\si{\celsius}, $221$\,\si{\bar})}

    % ---- Lineas de referencia (1 bar) ---------------------------------------
    \put(1.0, 2.0){\line(1, 0){7}}        % Linea horizontal p = 1 bar
    \put(8.1, 1.95){\tiny $p = 1$\,bar}

    % ---- Escala de presion --------------------------------------------------
    \put(0.6, 2.0){\small $1$}
    \put(0.3, 1.7){\small $6.1\times10^{-3}$}
    \put(0.6, 5.5){\small $221$}

    % ---- Escala de temperatura ----------------------------------------------
    \put(2.9, 0.25){\small $0.01$}
    \put(7.3, 0.25){\small $374$}

\end{picture}
\caption{Diagrama de fases esquem\'atico del agua (\ce{H2O}).
         Se muestran las tres curvas de coexistencia (s\'olido-l\'iquido,
         l\'iquido-vapor y s\'olido-vapor), el punto triple
         (\SI{0.01}{\celsius}, \SI{6.12e-3}{\bar}) y el punto cr\'itico
         (\SI{374}{\celsius}, \SI{221}{\bar}).
         N\'otese la pendiente \textit{negativa} de la curva de fusi\'on,
         caracter\'istica an\'omala del agua debida a que el hielo es menos
         denso que el agua l\'iquida.}
\label{fig:diagrama_fases_agua}
\end{figure}

% =============================================================================
% FIN DE figura_diagrama_fases_agua.tex
% =============================================================================

% =============================================================================

\begin{figure}[htbp]
\centering
\setlength{\unitlength}{1.2cm}
\begin{picture}(9, 7)

    % ---- Ejes ---------------------------------------------------------------
    \put(1.0, 0.5){\vector(1, 0){7.5}}   % Eje X (Temperatura)
    \put(1.0, 0.5){\vector(0, 1){6.0}}   % Eje Y (Presion)

    % Etiquetas de ejes
    \put(8.7, 0.25){\small $T$ / \si{\celsius}}
    \put(0.0, 6.6){\small $p$ / \si{\bar}}

    % ---- Lineas de las curvas de coexistencia --------------------------------

    % Curva solido-liquido (ligeramente inclinada a la izquierda — anomalia del agua)
    \put(3.5, 1.7){\line(-1, 4){0.9}}     % Pendiente negativa (dP/dT < 0)
    % Etiqueta
    \put(2.0, 4.5){\small\itshape Curva de}
    \put(2.0, 4.2){\small\itshape fusi\'on}

    % Curva liquido-vapor (desde punto triple hasta punto critico)
    \qbezier(3.5, 1.7)(5.5, 2.5)(7.5, 5.5)
    \put(5.8, 3.0){\small\itshape Curva de}
    \put(5.8, 2.7){\small\itshape vaporizaci\'on}

    % Curva solido-vapor (desde 0 hasta punto triple)
    \qbezier(1.2, 0.7)(2.0, 0.9)(3.5, 1.7)
    \put(1.5, 0.9){\small\itshape Curva de}
    \put(1.5, 0.65){\small\itshape sublimaci\'on}

    % ---- Regions (etiquetas de fases) ---------------------------------------
    \put(1.8, 3.5){\large S\'OLIDO}
    \put(4.2, 2.2){\large L\'IQUIDO}
    \put(4.8, 1.0){\large VAPOR}

    % ---- Punto triple -------------------------------------------------------
    \put(3.5, 1.7){\circle*{0.12}}
    \put(3.2, 1.45){\small Punto triple}
    \put(3.2, 1.22){\small ($0.01$\,\si{\celsius}, $6.12\times10^{-3}$\,\si{\bar})}

    % ---- Punto critico ------------------------------------------------------
    \put(7.5, 5.5){\circle*{0.12}}
    \put(7.0, 5.8){\small Punto cr\'itico}
    \put(7.0, 5.6){\small ($374$\,\si{\celsius}, $221$\,\si{\bar})}

    % ---- Lineas de referencia (1 bar) ---------------------------------------
    \put(1.0, 2.0){\line(1, 0){7}}        % Linea horizontal p = 1 bar
    \put(8.1, 1.95){\tiny $p = 1$\,bar}

    % ---- Escala de presion --------------------------------------------------
    \put(0.6, 2.0){\small $1$}
    \put(0.3, 1.7){\small $6.1\times10^{-3}$}
    \put(0.6, 5.5){\small $221$}

    % ---- Escala de temperatura ----------------------------------------------
    \put(2.9, 0.25){\small $0.01$}
    \put(7.3, 0.25){\small $374$}

\end{picture}
\caption{Diagrama de fases esquem\'atico del agua (\ce{H2O}).
         Se muestran las tres curvas de coexistencia (s\'olido-l\'iquido,
         l\'iquido-vapor y s\'olido-vapor), el punto triple
         (\SI{0.01}{\celsius}, \SI{6.12e-3}{\bar}) y el punto cr\'itico
         (\SI{374}{\celsius}, \SI{221}{\bar}).
         N\'otese la pendiente \textit{negativa} de la curva de fusi\'on,
         caracter\'istica an\'omala del agua debida a que el hielo es menos
         denso que el agua l\'iquida.}
\label{fig:diagrama_fases_agua}
\end{figure}

% =============================================================================
% FIN DE figura_diagrama_fases_agua.tex
% =============================================================================
